\begin{scanpage}{0231}{213}
\noindent it leads to inaccurate results in the case of odd resonances (especially for long conductors), and it does not give the displacement of the resonance peaks from the values $\lambda=2L/n$ $(n=1,3,5,\ldots)$ in the direction of longer wavelengths.

The results obtained by us are also approximate. However, our Equation (32.05) satisfies the reciprocity principle, and is applicable for any length vibrator. For very short vibrators $L\ll\lambda$, it changes into the asymptotic expression for the scattering characteristic of a dipole (see \S\ 33). For vibrators with a length of several wavelengths ($L\approx n\lambda$, $n=1,2,3,4$), Equation (32.05) gives satisfactory results. Calculations performed on its basis for radar reflection with normal irradiation agree with the results of Lindroth. Good agreement is also observed with the results of Leontovich and Levin for the integral scattering characteristic. With an increase of the vibrator's length, the precision of this equation increases, and in this way it is favorably distinguished from the equations proposed for the scattering characteristics by other authors.

Moreover, the divergence between the various approximation methods indicates the necessity of performing rather detailed calculations based on precise methods, for the purpose of evaluating the actual error of the approximation methods. Such calculations may be performed, for example, by means of the method discussed in References [88, 89] or [91].
\end{scanpage}
